% !TeX root = ../papers/nc-p3.tex
\section{全四辺でただ一つの単純な純関係の排除}
\label{sec:single-pure-branch-exclusion}
\begin{theorem}\label{thm:single-pure-branch-exclusion}
標準Koszul左右AS正則窓において、四つの $\mathbb P(R_i)$ とそれぞれのSegre多様体との
交叉スキームの非交和が長さ1になることはない。
\end{theorem}
\begin{proof}[中心の型と短い列の制限]
四辺が有限なら定理\ref{thm:finite-frobenius-low-opposite-as}による。
曲面以上の射影では定理\ref{thm:isolated-surface-fibers}を用いる。
非平面的な整曲面間の支配的対応では純関係を評価すると、片方の面が一つの超平面に入ることになり矛盾する。
平面列では純関係の一方の文字がその平面の法線であるが、その法線は正則正規である。
一次元の両側一点グラフは定理\ref{thm:curve-graph-pure-exclusion}による。
従って残るのは曲面射影のない分岐である。

唯一の純関係を $R_1$ に置く。命題\ref{prop:branch-local-pure-pencil}, \ref{prop:finite-pure-pencil-opposite}により、
普通の中心は全て $e=3$ で、例外の $e=2$ は順なら位置0、逆なら位置3にそれぞれ高々一つである。
少なくとも一方向に例外がある。例外中心は反対二次曲線に加えて、
ただ一つの例外方向でその曲線の一点を通る線を持つ。
両方向とも例外ならこの二中心は同じ二直線列の始終である。
実際、それらの三空間 $P,Q\subset R_1$ は同じ純関係を含み、その交叉はちょうど一次元である。
二次元なら、初回・終回直線 $L,M$ を消す $L^\perp\otimes M^\perp$ 内の
二つの行列により長さ2以上の純関係スキームを得る。
ゆえに $P+Q$ が五次元で $L\times M$ に残る双一次式は高々一本である。
零なら動く中心があり、階数1なら $L$ の一般点に同じ後点があり、動く位置が一つだけの列となる。
従ってその式は階数2で、初回直線から終回直線への移動を与える。

さらに、連続する位置 $i,i+1$ から二直線列が出て位置 $i+1$ の順中心が一つだけなら不可能である。
位置 $i$ からの初回直線上の順欠損はその一点に長さ2で集中する一方、直線の接触位数は1である。
命題\ref{prop:excess-contact-two-line-pure}から次の列の中央辺で長さ2以上の純関係が生じる。
また同じ位置から二本の $(1,3,1)$ 列が出て、反対位置の順中心が一つだけということもない。
両三次曲線は順欠損5を持ちその中心を通るが、そこでの一意な前点は
元の二中心の相異なる交わらない初回直線の両方には属せないからである。
\end{proof}

\begin{proof}[有限列挙と最後の二三次曲線]
各位置の順・逆中心数をそれぞれ $0,1,2$ から選ぶ。一方向の占有位相が一つだけなら、
そこからの列の初めの次数が $1,3,6$、第四次数が $10+b_2$ となり上界7に反する。
よって各方向に少なくとも二位相あり、普通の順中心が一つ以上ある。その中心の初回移動の位相では
初回以外は同じ反対三次曲線に着くから移動は二回以下であり、各列長は $m\le12$ である。

添付コードの有限列挙で使う条件は以下の通りである。
\begin{enumerate}
\item 順・逆中心の総数は等しい。各順中心から、適合する位相の相異なる逆中心まで
$2\le m\le12$ の列を割り当てる。両方向に例外があれば既述の二直線列を課す。
\item 全次数は正、初回・終回は1、同じ位相の次数和は7以下である。第二次数は $3-f_1$。
次位相に $n$ 中心があれば普通中心では全てが初回直線上にあり、例外中心でも
その外には高々一つしかない。外にある中心の後続直線は唯一の例外直線であって、二中心はこれを共有しない。
次位相が空なら $f_1=0$ である。自位相が二中心なら第二次数は相手の $e$ に等しい。
$m>2$ では第二次数は2以上である（1なら次数公式より第三次数は0以下）。
反対側でも同じ条件を課す。
\item 各移動がある中心の初回移動と同じ位相にありその初回自体でない場合、
行先成分を普通中心なら次数3、例外なら次数1か2に指定し、各中心につき同次数の行先を二度使わない。
この条件は反対曲線の各被約成分がその次数ごとに一本しかないことによる。逆向きも同じである。
\item 定理\ref{thm:pure-branch-degree-patterns}と同じ三階差の符号条件を課し、
前段の連続二直線列と同位相の二三次列の禁止を使う。
\end{enumerate}
例外位置に二中心があれば一番目を例外としてよい。整数計算で残るのは次の二行だけである。
\[
\begin{array}{c|c|c|c}
\text{順中心位相}&\text{逆中心位相}&\text{例外}&\text{二列の次数}\\\hline
0,2&0,2&\text{順 }0 &(1,3,1),(1,3,1)\\
1,3&1,3&\text{逆 }3 &(1,3,1),(1,3,1)
\end{array}
\]

最後は純関係の不在に依存せず定理\ref{thm:branch-pencil-resultant}の補完式を使える。
実際、二直線列がないので定理\ref{thm:ternary-center-cubic-scroll}が適用でき、
各初回直線と第二三次曲線への制限は $[s,t]A(t)$ で $A$ の一般階数は3である。
従って始点の関係空間を初回直線に制限した像は六次元である。異なる二つの一般 $t$ の
各々で $[-,t]$ と始点を消す全三空間の積を含むからである。逆方向でも同様である。
二つの直線が同じ位相に来るとき、それらは交わる。中央三次曲線が逆欠損5により逆中心を通り、
そのパラメータでの初回直線の点は逆中心の前点直線にも属する。両直線の一致は一動位置列を与えてしまう。
よって命題\ref{prop:branch-three-three-frame}の座標と $hh'B$ の表示が使え、
各奇位置の三座標制限は初回または終回直線を含むので前記の完全六空間を保つ。
両切断で $\det B\ne0$ となり、二つの補完式の共通四次根が両側中心を与える矛盾は同じである。
\end{proof}

\begin{corollary}\label{cor:four-two-line-impossible-with-pure}
純関係の個数や有限性を仮定しなくても、各位相一つずつの順・逆中心から四本の
二直線列ができる配置は標準正則窓にはない。
\end{corollary}
\begin{proof}
初回直線 $L_i$ 上にある順中心 $F_i$ は唯一の順欠損の台、終回直線 $M_i$ 上の $B_i$ も同様である。
$L_i$ から $M_{i+1}$ への同型で $B_{i+1}$ を引き戻せば $L_i\cap M_i$ の点となり、
両直線は一致しない。従って定理\ref{thm:final-pure-branch-frames}の四中心の基底、非零な $W$、
両切断の共通部分 $N$、その列空間が座標ベクトル2を含まないことは純関係があっても成立する。
同証明の最後で得た、初回直線を消す三関係のうち一つは、
したがって第四行だけに支持される純関係でなければならない。
巡回して全ての $R_i$ が $z_i$ に始まる純関係を持つ。ここで $z_i$ は
$L_i,M_i$ の張る平面を消す唯一の文字である。語順を逆にした同じ議論では、
全ての $R_{i-1}$ が同じ $z_i$ に終わる純関係を持つ。連続する純関係の禁止に反する。
\end{proof}
